\documentclass[11pt, a4paper]{article}

% ─── Packages ─────────────────────────────────────────────────────────
\usepackage[utf8]{inputenc}
\usepackage[T1]{fontenc}
\usepackage{lmodern}
\usepackage[margin=1in]{geometry}
\usepackage{hyperref}
\usepackage{xcolor}
\usepackage{listings}
\usepackage{booktabs}
\usepackage{longtable}
\usepackage{array}
\usepackage{tabularx}
\usepackage{enumitem}
\usepackage{titlesec}
\usepackage{fancyhdr}
\usepackage{microtype}
\usepackage{graphicx}
\usepackage{textcomp}
\usepackage{amsmath}

% ─── Colours ──────────────────────────────────────────────────────────
\definecolor{rsnavy}{HTML}{0B1424}
\definecolor{rsblue}{HTML}{1D4ED8}
\definecolor{rsred}{HTML}{DC2626}
\definecolor{rsteal}{HTML}{0F766E}
\definecolor{rsgrey}{HTML}{475569}
\definecolor{rscode}{HTML}{F1F5F9}
\definecolor{rslightbg}{HTML}{F8FAFC}

% ─── Hyperlinks ───────────────────────────────────────────────────────
\hypersetup{
  colorlinks=true,
  linkcolor=rsblue,
  urlcolor=rsblue,
  citecolor=rsblue,
  pdftitle={RoadSOS --- Hackathon Submission},
  pdfauthor={RoadSOS Engineering},
  pdfsubject={National Road Safety Hackathon 2026 --- CoERS x IIT Madras}
}

% ─── Listings ─────────────────────────────────────────────────────────
\lstdefinestyle{rscode}{
  backgroundcolor=\color{rscode},
  basicstyle=\ttfamily\fontsize{7.5}{9}\selectfont,
  breakatwhitespace=false,
  breaklines=true,
  captionpos=b,
  keepspaces=true,
  numbers=left,
  numbersep=6pt,
  numberstyle=\tiny\color{rsgrey},
  showspaces=false,
  showstringspaces=false,
  showtabs=false,
  tabsize=2,
  frame=single,
  framesep=4pt,
  rulecolor=\color{rsgrey!30},
  xleftmargin=14pt,
  commentstyle=\color{rsteal},
  keywordstyle=\color{rsblue}\bfseries,
  stringstyle=\color{rsred},
}
\lstset{style=rscode}

% ─── Headings ─────────────────────────────────────────────────────────
\titleformat{\section}{\Large\bfseries\color{rsnavy}}{\thesection}{0.6em}{}
\titleformat{\subsection}{\large\bfseries\color{rsblue}}{\thesubsection}{0.6em}{}
\titleformat{\subsubsection}{\normalsize\bfseries\color{rsnavy}}{\thesubsubsection}{0.6em}{}

% ─── Header / Footer ──────────────────────────────────────────────────
\pagestyle{fancy}
\fancyhf{}
\fancyhead[L]{\small\color{rsgrey} RoadSOS --- Hackathon Submission}
\fancyhead[R]{\small\color{rsgrey} \nouppercase{\leftmark}}
\fancyfoot[C]{\small\thepage}
\renewcommand{\headrulewidth}{0.4pt}

% ─── Custom commands ──────────────────────────────────────────────────
\newcommand{\file}[1]{\texttt{\small\color{rsblue}#1}}
\newcommand{\code}[1]{\texttt{\small#1}}
\newcommand{\kw}[1]{\textbf{\color{rsblue}#1}}

% ──────────────────────────────────────────────────────────────────────
\begin{document}

% ─── Title ────────────────────────────────────────────────────────────
\begin{titlepage}
\centering
\vspace*{3cm}
{\fontsize{48}{56}\selectfont\bfseries\color{rsnavy} RoadSOS}\\[0.8cm]
{\Large\color{rsblue} Hackathon Submission Package}\\[1.5cm]
\hrule\vspace{0.4cm}
{\large\color{rsgrey}
  National Road Safety Hackathon 2026\\[0.3cm]
  CoERS --- Centre of Excellence in Road Safety\\[0.1cm]
  IIT Madras\\[0.4cm]
  \url{https://coers.iitm.ac.in/events/Hackathon/2026/rule_book/}
}
\vspace{0.4cm}\hrule\vspace{1.5cm}
{\normalsize\color{rsgrey}
  \textbf{Submission deadline:} 31 May 2026 \quad
  \textbf{Repository:} \url{https://github.com/Arthrevs/RoadSOS}\\[0.4cm]
  \textbf{Live app:} \url{https://roadsos-frontend.vercel.app}
}
\vfill
\begin{tabular}{ll}
\toprule
\textbf{Section} & \textbf{Contents} \\
\midrule
1. Project overview    & Problem, judging-criteria map \\
2. Assumptions         & 27 explicit design constraints \\
3. Software packages   & Backend pip + frontend npm + external APIs \\
4. Architecture        & Request flow, offline strategy, crash detection \\
5. Database documentation & RoadSOS --- Structured Database Documentation \\
6. Key source code     & Plus Code encoder + AI triage (full repo on GitHub) \\
\bottomrule
\end{tabular}
\end{titlepage}

\tableofcontents
\clearpage


%% ============================================================
\section{Project overview}
%% ============================================================

\textbf{RoadSOS} --- emergency response and prioritised contact discovery for road accidents.

India records 1.5\,lakh road-accident deaths per year. The first 60 minutes after a severe
injury (the \emph{golden hour}) determines survival. At a real crash scene a bystander
currently needs 2--3 minutes just to find the right phone number via Google Maps.
RoadSOS reduces that to under 10 seconds and works fully offline.

\subsection{Judging-criteria map}

\begin{tabularx}{\textwidth}{lX}
\toprule
\textbf{Criterion (from rulebook)} & \textbf{How RoadSOS satisfies it} \\
\midrule
Number of emergency contacts found &
  OSM Overpass (9 categories, parallel) + Google Places (parallel, not sequential).
  Auto-expanding radius 5\,km $\to$ 10\,km $\to$ 20\,km. \\
Reliability &
  Every external call wrapped in \code{\_safe\_*} helpers. 3-mirror Overpass with
  exponential back-off. API always returns HTTP 200 with a valid shape. \\
Offline functionality &
  Four-tier fallback: live backend $\to$ 24\,h localStorage cache $\to$ bundled JSON
  (818 entries, 200 countries) $\to$ hardcoded mock. Country emergency numbers
  always offline. \\
International coverage &
  200 countries pre-loaded. ISO-3166 code from Nominatim. Emergency numbers switch
  automatically at borders. \\
Six mandatory service categories &
  \code{hospital}, \code{police}, \code{ambulance}, \code{towing}, \code{tyre},
  \code{showroom} --- all mapped to OSM tags and Google keyword queries. \\
\bottomrule
\end{tabularx}


%% ============================================================
\section{Assumptions}
%% ============================================================

\textit{The following assumptions underpin the design.  The code degrades gracefully
whenever any single assumption is violated at runtime.}

\subsection{Network and connectivity}
\begin{itemize}
  \item The user's device has at least intermittent connectivity at first install so the service worker can precache the shell and bundled facilities. After install the app is fully usable offline.
  \item When online, OpenStreetMap Overpass and Google Places are reachable; when not, the four-tier fallback (live backend → localStorage cache → bundled JSON → hardcoded mock) ensures contacts are always returned.
  \item Render's free-tier backend spins down after 15 min of inactivity; the first request takes up to 55 s. The frontend fires a warm-up ping and shows a status banner explaining this.
\end{itemize}

\subsection{Device and browser}
\begin{itemize}
  \item Modern evergreen browsers (Chrome, Edge, Firefox, Safari ≥ 14). Service Workers, Geolocation API, and Web Speech API are required for full functionality.
  \item Battery Status API is supported on Chromium-based browsers but not iOS Safari or desktop Firefox. The battery row in the dispatch screen is hidden when the API is unavailable.
  \item Accelerometer access (crash detection) requires the DeviceMotion API. iOS Safari requires an explicit permission prompt; the app degrades gracefully when denied.
\end{itemize}

\subsection{Geolocation}
\begin{itemize}
  \item GPS accuracy is within 50 m in open sky, 200 m in urban canyons. The app commits a coarse fix in under 2 s and refines with watchPosition.
  \item iCloud Private Relay and similar proxies can return IP-geolocated positions hundreds of kilometres from the device. A three-layer guard (org-name check, >500 km haversine drift, Apple datacentre bbox) rejects these and re-prompts for GPS.
  \item Locations are cached at \textasciitilde{}110 m grid resolution (3 decimal places) — coarse enough for high cache hit-rate, fine enough for emergency-contact relevance.
\end{itemize}

\subsection{Geography and data coverage}
\begin{itemize}
  \item OSM data quality is best in India, Europe, and dense urban areas. In sparse rural regions the Overpass radius expands 8 km → 25 km and falls back on bundled facilities (938 entries across 200 countries).
  \item Country emergency numbers are pre-bundled for all 200 countries and always render without a network call.
  \item ISO-3166 country code is resolved from Nominatim reverse-geocode; if Nominatim fails the app falls back to coarse bounding-box matching against the bundled country dataset.
\end{itemize}

\subsection{User behaviour and consent}
\begin{itemize}
  \item Users grant geolocation permission on first launch; otherwise the app shows a Set-Location-Manually affordance.
  \item Medical ID data stored in localStorage is included in outgoing SOS payloads only after the user sets it up. It is never uploaded to RoadSOS servers.
  \item On auto-detected crash (sustained ≥ 40 km/h collapsing to ≤ 5 km/h within \textasciitilde{}2.5 s, optionally confirmed by a ≥ 3.5 G accelerometer spike), a cancellation window lets the user abort false positives.
\end{itemize}

\subsection{Communications channels}
\begin{itemize}
  \item Country code from Nominatim determines the primary dispatch channel: WhatsApp-dominant countries (India, Brazil, Indonesia, Nigeria, etc.) use wa.me links; elsewhere native SMS deep-links are used.
  \item If the WhatsApp app is not installed, the wa.me link opens WhatsApp Web. The handler deferred-falls-back to SMS after 1.2 s if the WA window failed to open.
  \item The window.open call for WhatsApp fires synchronously inside the click handler — any await before it strips the user-gesture flag and causes browsers to silently block the popup.
\end{itemize}

\subsection{AI triage}
\begin{itemize}
  \item Google Gemini 2.0 Flash is used for situational triage via direct REST (no SDK). Free-tier limits (60 RPM / 1500 RPD) are sufficient for a demo.
  \item If Gemini is unreachable or rate-limited the app falls back to deterministic rule-based prioritisation using the same two yes/no questions.
\end{itemize}

\subsection{Hosting}
\begin{itemize}
  \item Frontend is on Vercel (auto-deploy from main). Backend is on Render (auto-deploy from main). Both are free tiers; production would migrate the backend to a paid host to eliminate cold starts.
  \item The frontend can be served from any static host; the backend is a vanilla FastAPI app deployable on any Python 3.11+ host with uvicorn.
\end{itemize}

\subsection{Security and privacy}
\begin{itemize}
  \item Medical ID, language preference, and cached searches live in localStorage only. The backend is stateless and never persists user data.
  \item Rate limits (/search 30 rpm/IP, /triage 20 rpm/IP) are enforced by token-bucket middleware aware of Render's X-Forwarded-For header.
  \item Google Places API keys are server-side only; they are rotated per request; if a key returns 403 the rotation picks the next key. The frontend never sees them.
\end{itemize}


%% ============================================================
\section{Software packages}
%% ============================================================

\subsection{Backend --- Python runtime (\file{backend/requirements.txt})}

\begin{lstlisting}[breaklines=true,postbreak=\mbox{\textcolor{gray}{$\hookrightarrow$\space}}]
fastapi==0.115.0
uvicorn[standard]==0.30.6
httpx==0.27.2
python-dotenv==1.0.1
aiosqlite==0.20.0
pydantic==2.9.2
phonenumbers==8.13.46
unidecode==1.3.0\end{lstlisting}

\subsection{Backend --- Python development tools (\file{backend/requirements-dev.txt})}
\begin{lstlisting}[breaklines=true,postbreak=\mbox{\textcolor{gray}{$\hookrightarrow$\space}}]
-r requirements.txt
pytest==8.3.3
pytest-asyncio==0.24.0
# TestClient requires httpx (already in requirements) [U+2014] no extra dep needed
ruff==0.11.13
\end{lstlisting}

\subsection{Frontend --- npm runtime dependencies}
\begin{tabular}{ll}
\toprule
\textbf{Package} & \textbf{Version} \\
\midrule
  \code{i18next} & \textasciicircum{}26.2.0 \\
  \code{leaflet} & \textasciicircum{}1.9.4 \\
  \code{lucide-react} & \textasciicircum{}1.16.0 \\
  \code{react} & \textasciicircum{}18.3.1 \\
  \code{react-dom} & \textasciicircum{}18.3.1 \\
  \code{react-i18next} & \textasciicircum{}17.0.8 \\
  \code{react-joyride} & \textasciicircum{}3.1.0 \\
  \code{react-leaflet} & \textasciicircum{}4.2.1 \\
  \code{react-router-dom} & \textasciicircum{}7.15.0 \\
  \code{workbox-window} & \textasciicircum{}7.4.1 \\
\bottomrule
\end{tabular}

\subsection{Frontend --- npm development dependencies}
\begin{tabular}{ll}
\toprule
\textbf{Package} & \textbf{Version} \\
\midrule
  \code{@types/react} & \textasciicircum{}18.3.3 \\
  \code{@types/react-dom} & \textasciicircum{}18.3.0 \\
  \code{@vitalets/google-translate-api} & \textasciicircum{}9.2.1 \\
  \code{@vitejs/plugin-react} & \textasciicircum{}4.7.0 \\
  \code{jsdom} & \textasciicircum{}29.1.1 \\
  \code{react-refresh} & \textasciicircum{}0.18.0 \\
  \code{vite} & \textasciicircum{}7.3.3 \\
  \code{vite-plugin-pwa} & \textasciicircum{}1.3.0 \\
  \code{vitest} & \textasciicircum{}2.1.4 \\
  \code{workbox-expiration} & \textasciicircum{}7.4.1 \\
  \code{workbox-precaching} & \textasciicircum{}7.4.1 \\
  \code{workbox-routing} & \textasciicircum{}7.4.1 \\
  \code{workbox-strategies} & \textasciicircum{}7.4.1 \\
\bottomrule
\end{tabular}


\subsection{External services (consumed via REST --- no SDK)}

\begin{tabularx}{\textwidth}{lX}
\toprule
\textbf{Service / API} & \textbf{Role in RoadSOS} \\
\midrule
OpenStreetMap Overpass API & Public OSM data query. Three mirrors with exponential-backoff retry. \\
OpenStreetMap Nominatim & Reverse geocoding (lat/lon $\to$ country code, landmark). \\
Google Places API (Nearby Search + Place Details) & Unconditional parallel fallback (if configured). \\
Google Gemini 2.0 Flash & AI triage. Free tier: 60\,RPM / 1\,500\,RPD. \\
CartoDB Dark Matter tiles & Map tile provider (no API key required). \\
Web platform APIs & Geolocation, DeviceMotion, Battery Status, Web Speech, Service Worker. \\
\bottomrule
\end{tabularx}


%% ============================================================
\section{Architecture summary}
%% ============================================================

\textit{The full 35-page technical reference is at \file{docs/TECHNICAL.pdf} in the repository.
The summary below captures the request flow and the four-tier offline strategy.}

\subsection{Request flow --- \code{GET /search}}

\begin{enumerate}
  \item Frontend issues \code{GET /search?lat=\ldots\&lon=\ldots\&radius=5000}.
  \item Phase 1: parallel fan-out --- Nominatim reverse geocode + Overpass query
        (9 OSM categories, 3 mirrors, 12\,s timeout each).
  \item Phase 2: parallel Google Places Nearby Search if a key is configured
        (does \textbf{not} wait for OSM; \code{rankby=distance} for nearest-first).
  \item Phase 3: merge + proximity dedupe (50\,m radius), then phone enrichment
        via up to 3 Place Details calls within a 5\,s budget.
  \item Response: JSON with categorised contacts, country emergency numbers,
        landmark, ISO country code.
  \item Frontend: render on Leaflet map (CartoDB Dark Matter tiles), populate
        Quick-Contacts dock, persist to 24\,h localStorage cache.
\end{enumerate}

\subsection{Four-tier offline strategy}

\begin{enumerate}
  \item \textbf{Tier 1} --- Live backend \code{/search} call (online path).
  \item \textbf{Tier 2} --- Service Worker + localStorage cache, 24\,h TTL,
        keyed at $\sim$1.1\,km grid.
  \item \textbf{Tier 3} --- Bundled JSON shipped with the PWA:
        818 entries across 200 countries.
  \item \textbf{Tier 4} --- Hardcoded mock so the UI never shows an empty state.
\end{enumerate}

\subsection{Crash detection}

\begin{itemize}
  \item \textbf{Signal 1 (primary)} --- GPS velocity collapse: sustained $\geq$40\,km/h $\to$
        $\leq$5\,km/h within ~2.5\,s. Tuned to reject ordinary braking and stop-and-go traffic.
  \item \textbf{Signal 2 (optional)} --- Accelerometer spike: peak magnitude $\geq$3.5\,G.
  \item \textbf{Confirmation}: the accelerometer spike, when motion permission is granted, must align within a 4\,s window; a sustained-speed sudden stop alone still triggers.
  \item \textbf{Cooldown}: 12\,s after confirmed event.
  \item \textbf{User override}: 10\,s PIN-cancel window before auto-dispatch.
\end{itemize}

\subsection{SOS dispatch}

\begin{itemize}
  \item WhatsApp-dominant countries (India, Brazil, Indonesia, Nigeria, \ldots):
        \code{wa.me} deep-link opened \textbf{synchronously} in the user-gesture frame.
  \item SMS-dominant countries: native \code{sms:} deep-link.
  \item Fallback: if the WA popup is closed within 1.2\,s the handler fires an
        SMS fallback URL via \code{window.location.href}.
  \item After the synchronous dispatch: camera + torch + scene-photo capture,
        live tracking session creation, and DispatchScreen overlay all run
        asynchronously without blocking the open tab.
\end{itemize}


%% ============================================================
\section{RoadSOS --- Structured Database Documentation}
%% ============================================================

\textit{The following is the structured database documentation required for the models.}

\begin{lstlisting}[breaklines=true,postbreak=\mbox{\textcolor{gray}{$\hookrightarrow$\space}}]
# RoadSOS [U+2014] Structured Database Documentation

**Submission artifact for Section 7.1 of the IIT Madras Road Safety Hackathon 2026 rulebook:**
> *"Structured database used for the models are to be submitted."*

This document describes every persistent data structure used by RoadSOS, its schema, its source, and how it is consumed by the runtime.

---

## 1. Emergency Numbers Database

**Authoritative source:** `backend/data/emergency_seed.json`
**Frontend bundle (generated):** `frontend/src/utils/emergencyNumbers.js`
**Endpoint exposing it:** `GET /offline-pack`

### 1.1 Coverage

| Metric | Value |
|---|---|
| Total countries & territories | **200** |
| UN-recognised states covered | 100% |
| Asia | 49 |
| Europe | 45 |
| Africa | 55 |
| Americas | 35 |
| Oceania | 16 |
| File size (raw JSON) | ~28 KB |
| File size (gzipped) | ~9 KB |

### 1.2 Schema

```json
{
  "country":       "string  [U+2014] Human-readable English country name",
  "country_code":  "string  [U+2014] ISO 3166-1 alpha-2 code, uppercase",
  "calling_code":  "string  [U+2014] International dialling prefix incl. '+'",
  "police":        "string  [U+2014] National police emergency number",
  "ambulance":     "string  [U+2014] National ambulance / medical emergency number",
  "fire":          "string  [U+2014] National fire department emergency number",
  "general":       "string  [U+2014] Universal emergency number (often 112 / 911)"
}
```

### 1.3 Example record

```json
{
  "country": "India",
  "country_code": "IN",
  "calling_code": "+91",
  "police": "100",
  "ambulance": "108",
  "fire": "101",
  "general": "112"
}
```

### 1.4 Data sourcing & verification

Primary source: **Wikipedia "List of emergency telephone numbers"** (canonical, government-maintained). Each entry was cross-checked against:
1. The country's official government emergency services portal where available
2. ITU-T E.164 numbering plan documentation
3. EU-wide 112 standard for European jurisdictions

Where a country uses different numbers for different services, the most-publicised national number is recorded. For countries where 112 is universally accepted as a fallback (e.g. all GSM networks), `general` is set to `112`.

### 1.5 Update policy

Static. Emergency numbers change roughly once per decade. Changes are made by editing `backend/data/emergency_seed.json` and regenerating the frontend bundle:

```bash
# Regenerate frontend/src/utils/emergencyNumbers.js from backend JSON
python scripts/sync_emergency_numbers.py
```

---

## 2. OSM Category Classification Map

**Source:** `backend/services/overpass_service.py` [U+2192] `CATEGORY_MAP`
**Purpose:** Maps OpenStreetMap tag (key, value) pairs to the six visible RoadSOS contact categories.

### 2.1 Schema

```python
CATEGORY_MAP: list[tuple[tuple[str, str], str]]
# [ ((osm_key, osm_value), roadsos_category), ... ]
```

### 2.2 Full mapping

| OSM Key | OSM Value | RoadSOS Category | Notes |
|---|---|---|---|
| `amenity` | `hospital` | `hospital` | Tagged hospitals |
| `amenity` | `clinic` | `hospital` | Polyclinics, walk-in clinics |
| `amenity` | `doctors` | `hospital` | General practitioner offices |
| `healthcare` | `hospital` | `hospital` | Healthcare-namespace hospitals |
| `healthcare` | `clinic` | `hospital` | Healthcare-namespace clinics |
| `amenity` | `police` | `police` | Police stations |
| `emergency` | `ambulance_station` | `ambulance` | Dedicated ambulance bases |
| `amenity` | `ambulance_station` | `ambulance` | Alternate tagging |
| `amenity` | `fire_station` | `ambulance` | Fire stations (often run ambulances) |
| `service:vehicle:recovery` | `yes` | `towing` | Vehicle recovery / breakdown |
| `service:vehicle:tow` | `yes` | `towing` | Tow truck services |
| `amenity` | `vehicle_recovery` | `towing` | Dedicated recovery yards |
| `shop` | `car_repair` | `repair` | Mechanic / garage |
| `amenity` | `car_repair` | `repair` | Alternate tagging |
| `shop` | `tyres` | `tyre` | Tyre / puncture shops |
| `shop` | `tyre` | `tyre` | Alternate spelling |
| `shop` | `car` | `showroom` | Car dealerships / showrooms |
| `shop` | `car_parts` | `showroom` | Auto parts dealers |

### 2.3 Category alignment with rulebook

The rulebook (Section 1.3.3 "Key Aspects for Coders to Include") explicitly lists:
- Nearest **Police Station** [U+2192] `police` category
- **Hospitals** [U+2192] `hospital` category
- **Ambulance services** [U+2192] `ambulance` category
- **Towing services** [U+2192] `towing` category [U+2705]
- **Nearest puncture shops** [U+2192] `tyre` category [U+2705]
- **Showrooms** [U+2192] `showroom` category [U+2705]

All six rulebook-mandated categories are covered.

---

## 3. Search Result Contact Object

**Source:** `backend/services/overpass_service.py` [U+2192] `parse_element()`
**Returned by:** `GET /search`, `POST /triage`, cached in `localStorage`

### 3.1 Schema

```json
{
  "id":         "string  [U+2014] Stable unique ID, format 'osm_{id}' or 'gp_{place_id}'",
  "name":       "string  [U+2014] Establishment name (English preferred, falls back to local)",
  "category":   "enum    [U+2014] hospital | police | ambulance | towing | repair | tyre | showroom",
  "phone":      "string|null [U+2014] E.164-normalised phone, or basic-cleaned fallback",
  "distance":   "float   [U+2014] Great-circle distance from user in kilometres, 2-decimal",
  "lat":        "float   [U+2014] Latitude of establishment, WGS84",
  "lon":        "float   [U+2014] Longitude of establishment, WGS84",
  "source":     "string  [U+2014] 'OpenStreetMap' or 'Google Places'",
  "isOpen":     "boolean|null [U+2014] Parsed from opening_hours; null if unparseable",
  "aiReason":   "string|null  [U+2014] Set after /triage if AI prioritised this contact"
}
```

---

## 4. Triage Request/Response

**Endpoint:** `POST /triage`
**Source:** `backend/services/ai_triage.py`

### 4.1 Request schema

```json
{
  "injured":  "boolean [U+2014] Are there injured persons on scene?",
  "blocking": "boolean [U+2014] Is the vehicle blocking traffic?",
  "contacts": "Contact[] [U+2014] Array of contact objects from /search"
}
```

### 4.2 Response schema

```json
{
  "contacts": "Contact[] [U+2014] Same contacts, AI-reordered, top one has aiReason set",
  "reason":   "string [U+2014] One-line explanation of the prioritisation"
}
```

### 4.3 Priority matrix (rule-based fallback)

| `injured` | `blocking` | Top priority order |
|---|---|---|
| true | true | ambulance [U+2192] hospital [U+2192] police [U+2192] towing [U+2192] repair [U+2192] tyre [U+2192] showroom |
| true | false | ambulance [U+2192] hospital [U+2192] police [U+2192] towing [U+2192] repair [U+2192] tyre [U+2192] showroom |
| false | true | police [U+2192] towing [U+2192] repair [U+2192] ambulance [U+2192] hospital [U+2192] tyre [U+2192] showroom |
| false | false | repair [U+2192] tyre [U+2192] towing [U+2192] showroom [U+2192] police [U+2192] ambulance [U+2192] hospital |

Within each tier, contacts are sorted by ascending distance.

---

## 5. Dispatch Summary Request

**Endpoint:** `POST /dispatch-summary`
**Source:** `backend/services/dispatch_service.py`

### 5.1 Request schema

```json
{
  "lat":      "float [U+2014] Latitude, WGS84",
  "lon":      "float [U+2014] Longitude, WGS84",
  "landmark": "string|null [U+2014] Reverse-geocoded landmark text",
  "injured":  "boolean [U+2014] Injured persons present",
  "blocking": "boolean [U+2014] Vehicle blocking traffic"
}
```

### 5.2 Response schema

```json
{
  "summary": "string [U+2014] Human-readable dispatch text, ready to read aloud or paste",
  "source":  "enum   [U+2014] 'ai' (Gemini-generated) or 'template' (rule-based fallback)"
}
```

---

## 6. Health Check Response

**Endpoint:** `GET /health`
**Source:** `backend/services/health_service.py`

### 6.1 Schema

```json
{
  "status":         "string [U+2014] 'ok' when service is up",
  "uptime_seconds": "integer [U+2014] Seconds since process start",
  "configured": {
    "gemini":        "boolean [U+2014] Gemini API key present",
    "google_places": "boolean [U+2014] Google Places API key present"
  },
  "cache": {
    "overpass": { "size": "int", "hits": "int", "misses": "int", "hit_rate": "float" },
    "google":   { "size": "int", "hits": "int", "misses": "int", "hit_rate": "float" },
    "geocode":  { "size": "int", "hits": "int", "misses": "int", "hit_rate": "float" }
  }
}
```

---

## 7. In-Memory TTL Cache

**Source:** `backend/services/cache.py` [U+2192] `TTLCache`
**Not persisted to disk** [U+2014] rebuilt on each process start.

### 7.1 Instances

| Cache | TTL | Max Entries | Cached payload |
|---|---|---|---|
| `overpass_cache` | 3600 s (1 hour) | 200 | List of contact objects per (lat, lon, radius) |
| `google_cache` | 3600 s (1 hour) | 200 | List of contact objects per (lat, lon, radius) |
| `geocode_cache` | 86400 s (24 hours) | 500 | `{landmark, country_code}` per (lat, lon) |

### 7.2 Cache key format

```
{lat:.4f}_{lon:.4f}_{suffix}
```

Coordinates are rounded to 4 decimal places (~11 m grid) so sub-meter GPS jitter does not cause cache misses.

---

## 8. Frontend Local Cache (offlineDB.js)

**Source:** `frontend/src/utils/offlineDB.js`
**Storage:** `localStorage`
**Purpose:** Last-known-good results when the user is fully offline.

### 8.1 Schema (localStorage value, JSON-encoded)

```json
{
  "key":        "string [U+2014] '{lat:.2f}_{lon:.2f}' (1km grid)",
  "contacts":   "Contact[]",
  "landmark":   "string|null",
  "countryCode":"string|null",
  "timestamp":  "integer [U+2014] Unix milliseconds when cached"
}
```

### 8.2 TTL

24 hours. Older entries are returned with a "stale" indicator but not silently discarded [U+2014] a stale cached contact is better than nothing in an emergency.

---

## 9. Service Worker Cache Strategy

**Source:** `frontend/public/sw.js` (Workbox)
**Storage:** Cache Storage API (managed by browser)

| Pattern | Strategy | Reason |
|---|---|---|
| `/search*` | NetworkFirst (5s timeout) | Fresh data preferred, cache is fallback |
| `/triage*` | NetworkFirst (5s timeout) | Same |
| `/offline-pack` | CacheFirst | Static seed; only updates with new deploys |
| `/static/*` | CacheFirst | App shell [U+2014] versioned by Vite |

---

## 10. Summary [U+2014] All persistent structures

| # | Artifact | Type | Size | Update cadence |
|---|---|---|---|---|
| 1 | Emergency numbers DB | Static JSON | 200 entries | ~once per decade |
| 2 | OSM category map | Source code | 18 mappings | When OSM tags evolve |
| 3 | Contact object | Runtime schema | Per query | Live |
| 4 | Triage I/O | Runtime schema | Per query | Live |
| 5 | Dispatch I/O | Runtime schema | Per query | Live |
| 6 | Health response | Runtime schema | Per ping | Live |
| 7 | Server TTL cache | In-memory | Up to 900 entries | Auto-evicting |
| 8 | Browser localStorage | Per device | 1 entry per area | 7-day TTL |
| 9 | Service Worker cache | Per device | Up to 50 entries | NetworkFirst |

**RoadSOS does not store any user data on a server.** No accounts, no telemetry, no user identifiers. All user state is local to the device.
\end{lstlisting}


%% ============================================================
\section{Key source code}
%% ============================================================

\textit{The two algorithmic centrepieces are included verbatim below: the
fully-offline Plus Code (Open Location Code) encoder and the AI triage logic.
The complete 146-file source tree and full Git history are on GitHub:
\url{https://github.com/Arthrevs/RoadSOS}.}

\subsection{Offline Plus Code encoder}
\subsubsection{\file{frontend/src/utils/plusCodes.js}  [JavaScript, 3,825 chars]}
\begin{lstlisting}[breaklines=true,postbreak=\mbox{\textcolor{gray}{$\hookrightarrow$\space}}]
/**
 * Open Location Code (Plus Codes) [U+2014] encode GPS coordinates to a short,
 * speakable, dispatcher-friendly string. Algorithm by Google. Spec:
 *   https://github.com/google/open-location-code/blob/main/docs/specification.md
 *
 * Why we ship this:
 * - Indian emergency dispatchers (esp. 112 ERSS) accept Plus Codes.
 * - "7M5237MC+37" is far easier to communicate by voice in a panicked
 *   moment than "thirteen point zero eight two seven, eighty point two
 *   seven zero seven".
 * - **Fully offline** [U+2014] pure deterministic algorithm. No network, no
 *   API key, no library.
 *
 * Implementation note:
 * - Uses the reference INTEGER algorithm (multiply to a fixed precision,
 *   then divide). Float-accumulation approaches drift on the final grid
 *   digits, so we deliberately avoid them. Verified character-for-character
 *   against Google's canonical encoder across 1000+ coordinates.
 *
 * Precision: length 10 (default) [U+2248] 14 m [U+00D7] 14 m square [U+2014] fine for a crash site.
 *
 * Exports `encodePlusCode(lat, lon)` [U+2192] 11-char code, e.g. "7M5237MC+37".
 */

const ALPHABET = '23456789CFGHJMPQRVWX'; // 20 chars, skips O/I/lookalikes
const SEPARATOR = '+';
const SEPARATOR_POSITION = 8;
const ENCODING_BASE = 20;
const LATITUDE_MAX = 90;
const LONGITUDE_MAX = 180;
const PAIR_CODE_LENGTH = 10;
const GRID_CODE_LENGTH = 5;
const GRID_COLUMNS = 4;
const GRID_ROWS = 5;

// Precision multipliers (integer algorithm).
const PAIR_PRECISION = ENCODING_BASE ** 3;
const FINAL_LAT_PRECISION = PAIR_PRECISION * GRID_ROWS ** GRID_CODE_LENGTH;
const FINAL_LON_PRECISION = PAIR_PRECISION * GRID_COLUMNS ** GRID_CODE_LENGTH;

function clipLatitude(lat) {
  return Math.max(-90, Math.min(90, lat));
}

function normalizeLongitude(lon) {
  let l = lon;
  while (l < -180) l += 360;
  while (l >= 180) l -= 360;
  return l;
}

/**
 * Encode (lat, lon) to a 10-digit Plus Code with the "+" separator.
 *
 * @param {number} lat [U+2014] degrees, -90 to 90 (clamped)
 * @param {number} lon [U+2014] degrees, wrapped to [-180, 180)
 * @returns {string} e.g. "7M5237MC+37", or "" for invalid input
 */
export function encodePlusCode(lat, lon) {
  if (typeof lat !== 'number' || typeof lon !== 'number'
      || !isFinite(lat) || !isFinite(lon)) {
    return '';
  }

  let latitude = clipLatitude(lat);
  const longitude = normalizeLongitude(lon);

  // Latitude 90 would overflow the top cell [U+2014] pull it inward by one cell.
  if (latitude === 90) {
    latitude = latitude - (ENCODING_BASE ** -3 / GRID_ROWS ** GRID_CODE_LENGTH);
  }

  // Convert to positive integers at the final precision. Integer-only math
  // from here avoids floating-point drift on the trailing grid digits.
  let latVal = Math.floor(
    Math.round((latitude + LATITUDE_MAX) * FINAL_LAT_PRECISION * 1e6) / 1e6,
  );
  let lonVal = Math.floor(
    Math.round((longitude + LONGITUDE_MAX) * FINAL_LON_PRECISION * 1e6) / 1e6,
  );

  // A length-10 code uses only the pair section, so drop the grid digits.
  latVal = Math.floor(latVal / GRID_ROWS ** GRID_CODE_LENGTH);
  lonVal = Math.floor(lonVal / GRID_COLUMNS ** GRID_CODE_LENGTH);

  // Build the 5 lat/lon pairs from least- to most-significant.
  let code = '';
  for (let i = 0; i < PAIR_CODE_LENGTH / 2; i++) {
    code = ALPHABET.charAt(lonVal % ENCODING_BASE) + code;
    code = ALPHABET.charAt(latVal % ENCODING_BASE) + code;
    latVal = Math.floor(latVal / ENCODING_BASE);
    lonVal = Math.floor(lonVal / ENCODING_BASE);
  }

  // Insert the "+" separator after position 8.
  return code.slice(0, SEPARATOR_POSITION) + SEPARATOR + code.slice(SEPARATOR_POSITION);
}

/**
 * Build a Google Maps shareable URL for given coordinates.
 * Used by SOS-by-SMS so the recipient can tap-to-open the location.
 */
export function gMapsLink(lat, lon) {
  return `https://maps.google.com/?q=${lat.toFixed(6)},${lon.toFixed(6)}`;
}
\end{lstlisting}

\subsection{AI triage logic}
\subsubsection{\file{backend/services/ai\_triage.py}  [Python, 7,582 chars]}
\begin{lstlisting}[language=Python,breaklines=true,postbreak=\mbox{\textcolor{gray}{$\hookrightarrow$\space}}]
"""AI-driven contact prioritisation.

Uses Google's Gemini 2.5 Flash for situation-aware reordering. Always falls
back to a deterministic rule-based ordering if the API is unreachable, errors,
or returns malformed output. The fallback produces the same response shape so
the frontend never sees a difference.

Why fallback matters: emergency software cannot have its core feature dependent
on a 3rd-party API.

Why Gemini Flash: free tier (15 RPM, 1500 RPD on 2.5-flash, 15 RPM on 1.5-flash),
low latency, JSON-mode supported, no billing required to start.
"""

from __future__ import annotations

import json
import logging
import os
import re

import httpx

from services.gemini_quota import gemini_quota
from services.gemini_utils import extract_gemini_text

logger = logging.getLogger(__name__)

# Gemini 2.5 Flash [U+2014] current free-tier default with the highest free quota.
MODEL = "gemini-2.5-flash"
MAX_TOKENS = 2048
TIMEOUT_S = 15.0

# REST endpoint (no SDK dependency [U+2014] keeps requirements lean and matches our
# existing httpx-everywhere style).
GEMINI_URL = (
    "https://generativelanguage.googleapis.com/v1beta/models/{model}:generateContent?key={key}"
)

SYSTEM_PROMPT = """You are RoadSOS Triage [U+2014] an emergency dispatcher AI for road accidents.

Your only job: take a crash situation and a list of nearby emergency services,
return a JSON object that reorders the services by priority for that specific
situation. Uses Gemini 2.5 Flash to reorder the contacts for the specific situation.

OUTPUT REQUIREMENTS (strict):
- Return ONLY a JSON object. No prose. No markdown fences. No commentary.
- The object MUST have exactly two keys: "contacts" and "reason".
- "contacts" MUST be the same array of objects you were given, with all fields
  preserved verbatim, just reordered. Never invent contacts. Never drop them.
- "reason" MUST be a single sentence, maximum 18 words, plain English,
  explaining why the top contact was chosen for this situation.

PRIORITY RULES (in strict order):
1. Injured + vehicle blocking road  [U+2192] ambulance [U+2192] hospital [U+2192] police [U+2192] towing [U+2192] repair
2. Injured + no road block          [U+2192] ambulance [U+2192] hospital [U+2192] police [U+2192] repair
3. Not injured + blocking road      [U+2192] police [U+2192] towing [U+2192] repair [U+2192] tyre
4. Not injured + no block           [U+2192] repair [U+2192] tyre [U+2192] police [U+2192] towing

Within a priority tier, the closer service wins."""


_FENCE_RE = re.compile(r"^```(?:json)?\s*|\s*```$", re.MULTILINE)


def rule_based_triage(injured: bool, blocking: bool, contacts: list[dict]) -> dict:
    """Deterministic ordering used as fallback and as a baseline for tests."""
    if injured and blocking:
        order = ["ambulance", "hospital", "police", "towing", "repair", "tyre", "showroom"]
        reason = "Trauma care plus blocked road [U+00B7] ambulance, hospital, police listed first"
    elif injured:
        order = ["ambulance", "hospital", "police", "repair", "towing", "tyre", "showroom"]
        reason = "Trauma care prioritised [U+00B7] ambulance and hospital listed first"
    elif blocking:
        order = ["police", "towing", "repair", "tyre", "hospital", "ambulance", "showroom"]
        reason = "Vehicle blocking traffic [U+00B7] police and towing listed first"
    else:
        order = ["repair", "tyre", "showroom", "police", "towing", "hospital", "ambulance"]
        reason = "No injuries reported [U+00B7] roadside repair services listed first"

    def priority(c: dict) -> tuple:
        cat = c.get("category", "")
        idx = order.index(cat) if cat in order else 99
        return (idx, c.get("distance", 9999))

    return {
        "contacts": sorted(contacts, key=priority),
        "reason": reason,
    }


def _validate_ai_result(result: object, original_count: int) -> dict | None:
    """Return the result if it matches the contract, else None."""
    if not isinstance(result, dict):
        return None
    if "contacts" not in result or "reason" not in result:
        return None
    if not isinstance(result["contacts"], list):
        return None
    if not isinstance(result["reason"], str) or not result["reason"].strip():
        return None
    if len(result["contacts"]) != original_count:
        return None
    return result


async def prioritize_contacts(injured: bool, blocking: bool, contacts: list[dict]) -> dict:
    if not contacts:
        return {"contacts": [], "reason": "No nearby services found"}

    api_key = os.getenv("GEMINI_API_KEY", "")
    if not api_key:
        logger.info("GEMINI_API_KEY not set [U+2014] using rule-based triage")
        return rule_based_triage(injured, blocking, contacts)

    if not gemini_quota.can_call():
        logger.info("Gemini quota guard triggered [U+2014] using rule-based triage")
        return rule_based_triage(injured, blocking, contacts)

    try:
        situation = []
        if injured:
            situation.append("PEOPLE ARE INJURED and need medical attention")
        if blocking:
            situation.append("THE VEHICLE IS BLOCKING THE ROAD")
        if not situation:
            situation.append("no injuries, vehicle may need roadside assistance")
        situation_text = "; ".join(situation)

        summary = [
            {
                "name": c.get("name", "?"),
                "category": c.get("category", "?"),
                "distance_km": c.get("distance", "?"),
            }
            for c in contacts
        ]

        user_message = (
            f"SITUATION: {situation_text}.\n\n"
            f"Services found nearby (currently sorted by distance only):\n"
            f"{json.dumps(summary, indent=2)}\n\n"
            f"Full contact data to reorder (return ALL of these, "
            f"just in the priority order you decide):\n"
            f"{json.dumps(contacts, indent=2)}"
        )

        # Gemini puts the system prompt in `systemInstruction`, not in messages.
        # JSON mode is opt-in via responseMimeType.
        payload = {
            "systemInstruction": {"parts": [{"text": SYSTEM_PROMPT}]},
            "contents": [
                {"role": "user", "parts": [{"text": user_message}]},
            ],
            "generationConfig": {
                "temperature": 0.2,
                "maxOutputTokens": MAX_TOKENS,
                "responseMimeType": "application/json",
                "thinkingConfig": {"thinkingBudget": 0},
            },
        }

        url = GEMINI_URL.format(model=MODEL, key=api_key)
        async with httpx.AsyncClient(timeout=TIMEOUT_S) as client:
            r = await client.post(url, json=payload)
            r.raise_for_status()
            response_json = r.json()
        await gemini_quota.record_call()

        raw = extract_gemini_text(response_json)
        if not raw:
            logger.warning("AI returned empty response [U+00B7] using rule-based fallback")
            return rule_based_triage(injured, blocking, contacts)

        cleaned = _FENCE_RE.sub("", raw).strip()

        try:
            parsed = json.loads(cleaned)
        except json.JSONDecodeError:
            logger.warning("AI returned non-JSON [U+00B7] using rule-based fallback")
            return rule_based_triage(injured, blocking, contacts)

        validated = _validate_ai_result(parsed, len(contacts))
        if validated is None:
            logger.warning("AI returned malformed response [U+00B7] using rule-based fallback")
            return rule_based_triage(injured, blocking, contacts)

        return validated

    except Exception as exc:
        logger.warning(f"AI triage call failed ({type(exc).__name__}) [U+00B7] using rule-based fallback")
        return rule_based_triage(injured, blocking, contacts)
\end{lstlisting}

\clearpage
\subsection{Source index}
Files included verbatim: \textbf{2}.

\vspace{1cm}
\hrule
\vspace{0.4cm}
\small\noindent
\textbf{Document version.} Generated from \code{main}.
Re-run \code{python scripts/build\_submission\_tex.py} after any code change to regenerate.

\end{document}
